\documentclass[conference]{IEEEtran}
\IEEEoverridecommandlockouts
\usepackage{cite}
\usepackage{amsmath,amssymb}
\usepackage{graphicx}
\usepackage{booktabs}
\usepackage{url}
\usepackage[hidelinks]{hyperref}

\begin{document}

\title{Ticker Extraction from Financial Influencer Videos}

\author{\IEEEauthorblockN{Sri Veerabhadra Karthik Mulugu}
\IEEEauthorblockA{\textit{Incoming Ph.D. Student, Computer Science} \\
\textit{The University of Texas at San Antonio}\\
San Antonio, TX, USA \\
karthikmulugu14@gmail.com}}

\maketitle

\begin{abstract}
The VideoConviction benchmark poses stock ticker extraction from financial
influencer videos as a generation task: a model reads a transcript segment and
emits a symbol. Before building an extractor I measured whether the symbol is
present to be read, and found it absent from 47.9\% of annotated segments. The
absence is not uniform. It concentrates on the largest companies: the symbol is
spoken in 60.5\% of segments discussing sub-\$300M firms but only 28.1\% of
those discussing firms above \$200B ($\rho=-0.215$, $p=1.5\times10^{-4}$). TSLA
is annotated in 26 segments and the string ``TSLA'' occurs in none of them. A
method that searches for the symbol is therefore not merely imperfect but
systematically biased against the recommendations carrying the most capital. I
reformulate the task as retrieval followed by ranking. A video-level stage
harvests candidates from the description, title, full transcript, dense
retrieval over company names, and optical character recognition of sampled
frames, reducing a 10{,}550-symbol universe to a median of 93 candidates while
retaining the correct answer 92.6\% of the time. A segment-level LambdaMART
ranker then selects among them using 37 features. Evaluated out-of-fold with
folds grouped by video and scores averaged over three random seeds, the system
attains F1 80.06, exceeding every text-only model reported in the benchmark
including Qwen2.5-72B at 79.20, while using no large language model and no paid
inference. I report two implementation defects in full, because both were cases
where a correction in one component opened a failure in another, and both were
invisible in aggregate metrics.
\end{abstract}

\begin{IEEEkeywords}
information extraction, learning to rank, multimodal retrieval, optical
character recognition, financial NLP, selective prediction
\end{IEEEkeywords}

\section{Introduction}

Financial influencers, colloquially ``finfluencers,'' publish stock
recommendations to large audiences on video platforms. Studying their influence
at scale requires first identifying automatically which stock each
recommendation concerns.

\subsection{What the Benchmark Provides}

VideoConviction~\cite{videoconviction} supports this with more than 6{,}000
expert annotations produced over 457 hours of human labelling. Each annotation
marks a timestamped segment of a video containing exactly one recommendation,
recording the ticker symbol, the recommended action, and a conviction score
reflecting how strongly the presenter appears to believe the claim. The authors
define three tasks of increasing difficulty: extracting the ticker (T), the
ticker and action (TA), and the ticker, action and conviction (TAC). Task T is
first, because the wrong symbol makes the action and conviction wrong
too.

\subsection{What the Benchmark Does}

The authors evaluate 21 language and multimodal models on all three tasks. Their
inference procedure is stated in a single sentence: a single prompt instructs
the model to extract ticker, action and conviction simultaneously, zero-shot.
No symbol registry is supplied, and the annotation field contains the correct
ticker with no corresponding company name. Models therefore \emph{generate} a
symbol from parametric knowledge.

The strongest text-only system reaches F1 79.20 and the strongest multimodal
system 86.23. The authors attribute the multimodal advantage to on-screen charts
displaying ticker names, and separately observe that models emit plausible but
non-existent symbols, giving \texttt{APPL} for Apple as their example.

\subsection{What I Did Differently, and Why}

Rather than start by building an extractor, I first measured whether the
answer is in the input at all. Often it is not, and the pattern of what is
missing shaped every choice that followed.

\begin{enumerate}
\item \textbf{A measurement that changes the task
(Section~\ref{sec:char}).} The ticker symbol is absent from 47.9\% of segments,
and absence correlates negatively with company size, so symbol matching fails
hardest on the economically largest recommendations. I quantify this with a rank
correlation, a bootstrap confidence interval and a non-parametric test.

\item \textbf{Splitting the task into two stages
(Sections~\ref{sec:rationale}--\ref{sec:rank}),} adopted because a simpler approach was
measured and found worse, not by assumption. Each stage is measured separately, so the
binding constraint on system performance is explicit and attributable.

\item \textbf{Choosing from a list instead of generating (Section~\ref{sec:universe}).}
Choosing from a constructed 10{,}550-symbol registry eliminates the
hallucination failure the benchmark names, by construction rather than by
prompting. It also removes the requirement that a model have memorised the
symbol of a company delisted years ago, which matters because roughly a quarter
of the correct answers are in that category.

\item \textbf{A cheap way to use the video
(Section~\ref{sec:ocr}).} If the multimodal advantage arises because symbols are
\emph{written} on screen, what is needed is reading text, not
understanding images. Frame OCR recovers the symbol in 41.9\% of segments and
lifts the system from F1 78.02 to 80.06.
\end{enumerate}

I additionally report two defects found during development
(Sections~\ref{sec:idfbug} and~\ref{sec:compound}). I include them because the
first failed on precisely the population my own analysis had identified as
critical, and because neither was visible in an aggregate score.

\section{Data and Preprocessing}

\subsection{Structure}

Each row of VideoConviction corresponds to one segment of one video concerning
one stock; a video covering five stocks yields five rows. The fields I use are
the segment transcript, the full video transcript, the video title and
description, the segment start and end timestamps, and the source URL.

Transcripts are produced by automatic speech recognition (ASR). This is
consequential throughout: company names are routinely corrupted, and
orthographic conventions do not survive transcription.

\subsection{Filtering}

Of 760 annotated rows, 687 carry \texttt{is\_rec\_present = Yes}. The remaining
73 are segments in which a stock is discussed but no recommendation is made.
Since Task T asks which stock is being \emph{recommended}, these have no target,
and I exclude them. All results below are over the 687. No segment in that set
has an empty transcript, so no further filtering is applied.

\subsection{Metric}

I report F1 against the human annotations, following the benchmark. I
additionally report coverage and precision separately, because a system
permitted to decline is more useful downstream than one forced to commit, and
that distinction is invisible in a single F1 figure.

\section{Characterising the Task}
\label{sec:char}

\subsection{Is the Answer Present?}

For each of the 687 segments I asked whether the correct symbol occurs in the
transcript in any surface form. I tested three, chosen by reading transcripts
rather than taken from prior work: a bare token (\texttt{MITK}), a cashtag
(\texttt{\$MITK}), and a dictated spelling (\texttt{M-I-T-K}), the last being
common when a presenter reads out an unfamiliar symbol.

\begin{table}[htbp]
\caption{Surface realisation of the correct ticker in the segment transcript
($n = 687$). Forms are not mutually exclusive.}
\label{tab:surface}
\centering
\begin{tabular}{lrr}
\toprule
Surface form & Segments & Share \\
\midrule
Cashtag \texttt{\$TICK} & 0 & 0.0\% \\
Bare token \texttt{TICK} & 333 & 48.5\% \\
Dictated \texttt{T-I-C-K} & 38 & 5.5\% \\
\midrule
Any literal form & 358 & 52.1\% \\
\textbf{No literal form} & \textbf{329} & \textbf{47.9\%} \\
\bottomrule
\end{tabular}
\end{table}

Two observations follow from Table~\ref{tab:surface}.

First, no cashtag occurs in any transcript. The \texttt{\$TICK} convention is
the standard marker in financial social-media text and I expected it to be the
primary signal. It is orthographic, never spoken, and therefore never produced
by ASR. I note the zero explicitly because anyone importing methods from
StockTwits or Twitter corpora will build on a convention that does not survive
this medium. The convention does appear in the \emph{typed} fields, 37 times in
descriptions and 7 in titles, which is consistent with the same explanation.

Second, and more consequentially, nearly half of all segments contain no form of
the target string. For those segments the answer is not in the input, and no
improvement in matching can recover it.

\subsection{Absence Is Not Random}
\label{sec:inversion}

I then asked whether the 329 segments lacking a symbol differ systematically
from those containing one. Binning by the market capitalisation of the correct
ticker yields Table~\ref{tab:cap}.

\begin{table}[htbp]
\caption{Symbol visibility by market capitalisation}
\label{tab:cap}
\centering
\begin{tabular}{lrr}
\toprule
Capitalisation band & Segments & Symbol spoken \\
\midrule
$<$ \$300M & 38 & 60.5\% \\
\$300M--\$2B & 40 & 55.0\% \\
\$2B--\$10B & 76 & 64.5\% \\
\$10B--\$200B & 211 & 60.7\% \\
$>$ \$200B & 167 & \textbf{28.1\%} \\
\bottomrule
\end{tabular}
\end{table}

I tested the pattern three ways rather than reporting the table as a trend,
and chose each test for a reason.

\emph{Spearman's rank correlation} ($\rho = -0.215$, $p = 1.5\times10^{-4}$,
$n = 307$ tickers) asks whether two quantities move together using their rank
order rather than their raw values. Rank order is right here because market
capitalisation spans five orders of magnitude, so one trillion-dollar company
would dominate an ordinary correlation. It is computed per ticker, not per
segment, so a single heavily covered stock cannot drive the result.

\emph{A bootstrap confidence interval} on the gap between the bottom and top
capitalisation terciles gives $19.4$ points, 95\% interval $[7.1, 31.4]$. It is
obtained by drawing 5{,}000 resamples of each group with replacement,
recomputing the gap each time, and reading the middle 95\% of those values. It
asks how far the answer would move on a different sample of the same size, and
the interval excluding zero rules out the effect being an artefact of which
companies happen to appear here.

\emph{A one-sided Mann--Whitney test} ($p = 5.1\times10^{-4}$) checks the same
comparison without assuming the data is normally distributed. Visibility per
ticker is bounded in $[0,1]$ and clusters at the endpoints, so a $t$-test would
not be appropriate.

The relationship is better described as a threshold than a gradient. Visibility
is approximately flat between 55\% and 67\% across four bands and collapses only
in the mega-capitalisation band. The explanation is simple: a speaker says the symbol
when the audience cannot be assumed to know the company, and that changes
sharply at the point of household recognition rather than smoothly with size.

The individual cases are stark. TSLA is annotated in 26 segments and the string
``TSLA'' appears in none of them; AAPL in 10, likewise none. Conversely a
\$225M firm has its symbol dictated character by character: \emph{``Veritone
Inc., ticker V-E-R-I.''}

Because capitalisation figures are present-day while the videos span 2019--2023,
companies that have since delisted are excluded from this analysis (155 of 687
segments). Delisted firms are disproportionately small, so their exclusion
removes cases that would strengthen the correlation. The measured effect is
conservative.

\subsection{Why This Determines the Architecture}

The problem is not just that symbol matching misses cases. It is that it
misses them in a pattern. A system built on matching the symbol
string performs worst on precisely the recommendations with the largest market
impact, and an aggregate score conceals this. I therefore treat the symbol as
one evidence source among several rather than as the object of extraction.

\section{Design Rationale}
\label{sec:rationale}

\subsection{A Test That Made Things Worse}

Given that the symbol is often absent from the segment, the natural remedy is to
supply more text. The dataset includes each video's description, in which I
measured the correct symbol to appear 35.1\% of the time. Supplying it to a flat
extractor produced Table~\ref{tab:negresult}.

\begin{table}[htbp]
\caption{Effect of adding the video description to a flat extractor}
\label{tab:negresult}
\centering
\begin{tabular}{lrrr}
\toprule
Input & Coverage & Precision & Recall \\
\midrule
Segment transcript only & 61.7\% & \textbf{71.2\%} & 44.0\% \\
+ video description & 86.5\% & \textbf{51.0\%} & 44.1\% \\
\bottomrule
\end{tabular}
\end{table}

Recall moved by 0.1 points while precision fell by 20. Supplying additional
genuinely relevant information made the system substantially worse.

\subsection{Diagnosis}

Inspecting the induced errors explains it. A finfluencer's description
enumerates every stock covered by the video, typically as a timestamped index.
The field therefore identifies the candidate set for the \emph{video} and
carries no information about which candidate a given \emph{segment} concerns.
The flat extractor was selecting the first symbol it encountered, usually a
different stock from the same video.

The description is not a feature. It is a candidate list.

\subsection{The Resulting Decomposition}

That generalises. A signal that narrows the field without choosing between the
options belongs in a retrieval stage, not in a scoring function. I therefore decompose Task T into a video-level stage that
must not omit the answer and a segment-level stage that must not guess. This
retrieve-and-rank pattern is standard in document ranking, where an inexpensive
high-recall pass precedes an expensive precise one~\cite{luan}; its
applicability here follows from how these videos are structured rather than
from convention.

\section{Stage 1: Candidate Generation}

\subsection{Constructing a Candidate Universe}
\label{sec:universe}

The benchmark supplies no symbol registry, so I built one.

This is a substantive design difference rather than plumbing. The evaluated
systems generate a symbol from parametric knowledge, which is the direct cause
of the failure the benchmark names: a generated string such as \texttt{APPL} is
not checked against anything. Selecting from a registry removes that failure by
construction, since an absent symbol cannot be emitted.

A second consequence matters more here. A generative model must have memorised
that Zynex trades as \texttt{ZYXI}. A selective model need only recognise that
the segment concerns Zynex, because the symbol lookup happens deterministically
upstream. Approximately 24\% of the correct answers belong to companies delisted
before or shortly after the evaluated models' training cutoffs, where that
memorisation is least likely to be present.

I assembled the registry from three sources, in order of authority.

\emph{The SEC.} The Securities and Exchange Commission publishes
\texttt{company\_tickers.json}, listing every currently registered company:
10{,}434 symbols, authoritative and free. It covers 301 of the 441
correct answers.

\emph{yfinance.} The SEC file is \emph{current}, and the videos span
2019--2023. Companies acquired, renamed or bankrupted since are absent: Twitter,
Slack, Xilinx, FireEye, Activision Blizzard. I queried Yahoo Finance for each of the
140 correct symbols the SEC file missed, because its coverage is historical
rather than current. It resolved 34 of them.

\emph{Manual entry.} The remaining 106 were resolved by hand where a name
could be established at all, which it could for 82. The count is not a budget I
chose; it is what was left after the two automatic sources, and every one of
them is a correct answer the system could otherwise never retrieve by name.

The resulting registry has 10{,}550 entries. Candidates are always drawn from
all of it, never from the 441 symbols that are correct answers; restricting the
candidate set to the label set would inflate precision by a large and
uninterpretable margin.

Twenty-four of the 441 correct answers resolve to no name from any source and
are left blank rather than guessed, since a wrong name would actively steer the
matcher toward the wrong company. Four are not symbols at all but annotation
errors; the remainder are delisted over-the-counter and foreign shells. Those
segments can therefore only be reached through the symbol itself, and they
account for part of the gap between stage-1 recall and 100\%.

\subsection{Evidence Sources}

For each video I pool evidence from the segment transcript, the full transcript,
the title and the description. The detection rules, listed in Table~\ref{tab:rules}, were written by
inspecting transcripts; none is taken from the benchmark, which describes no
detection procedure.

\begin{table}[htbp]
\caption{Detection rules and how often each fires}
\label{tab:rules}
\centering
\begin{tabular}{lr}
\toprule
Rule & Fires on \\
\midrule
Cue phrase (``ticker VERI'', ``trades under'') & 34.4\% of segments \\
Dictated spelling \texttt{V-E-R-I} & 5.4\% \\
Cashtag \texttt{\$VERI} & 0\% spoken, 44 typed \\
Bare capitalised token in registry & frequent, low precision \\
Company name via inverted index & --- \\
Dense retrieval over name embeddings & --- \\
Symbol read by frame OCR & 41.9\% \\
\bottomrule
\end{tabular}
\end{table}

Bare capitalised tokens are filtered against a list of symbols that are also
ordinary English words (\texttt{ALL}, \texttt{NOW}, \texttt{WELL}), without
which the rule produces almost only false positives. Measured in isolation, the
cue-phrase rule combined with a dictated spelling is the most precise at 96.6\%,
while the bare-token rule reaches only 50.5\%.

\subsection{Weighting Name Evidence}

Name matching requires a notion of how much a given word constrains the identity
of a company. I use inverse document frequency (IDF) over two distinct reference
collections, because either alone is insufficient.

\emph{Gazetteer IDF} measures how few company names contain a token. It
correctly assigns zero weight to \emph{Holdings} and high weight to
\emph{Veritone}. It cannot distinguish a distinctive proper noun from an
ordinary English word that happens to be rare in company names, so it rates
\emph{Next} and \emph{Today} highly. An early version of this system
consequently matched the word ``today'' to a company.

\emph{Corpus IDF} measures how few \emph{transcripts} contain a token, learned
from the dataset itself. This suppresses ordinary vocabulary without any
hand-authored stoplist, which matters because a stoplist I wrote would encode
my assumptions about what is ordinary rather than measure it.

Both are computed as $\log(N / n_t)$, where $N$ is the number of documents in
the reference collection and $n_t$ the number containing token $t$. For
gazetteer IDF the documents are the 10{,}550 company names; for corpus IDF they
are the 687 transcripts. The two are multiplied, so a token must be rare in both
senses to carry weight: distinctive as a company name \emph{and} unusual in
finance speech.

\begin{table}[htbp]
\caption{The two IDFs and their product on representative tokens}
\label{tab:idf}
\centering
\begin{tabular}{lrrr}
\toprule
Token & Gazetteer IDF & Corpus weight & Product \\
\midrule
veritone & 9.26 & 1.00 & 9.26 \\
skin & 9.26 & 0.78 & 7.20 \\
\textbf{tesla} & \textbf{9.26} & \textbf{0.28} & \textbf{2.56} \\
today & 8.57 & 0.18 & 1.52 \\
holdings & 0.00 & 0.53 & 0.00 \\
\bottomrule
\end{tabular}
\end{table}

\subsection{Defect 1: Corpus-IDF Suppression of Large Caps}
\label{sec:idfbug}

The \emph{tesla} row of Table~\ref{tab:idf} is the defect. Its gazetteer IDF
equals that of \emph{veritone}, but its corpus weight is 0.28 because Tesla is
discussed in a large fraction of transcripts. The product collapses.

I had applied the product at both stages. At the retrieval stage this removed
TSLA (17 segments), AAPL (7), NVDA (6) and AMZN (6) from the candidate pool
entirely, before any ranking occurred.

The structure of the error is worth stating plainly. Corpus IDF is a frequency
penalty; the most frequently discussed companies are the mega caps; and the mega
caps are precisely the population identified in Section~\ref{sec:inversion} as
the one where the symbol is least often spoken and name evidence therefore most
needed. A correction I introduced for precision had opened a recall failure in
the exact population my own analysis had flagged as critical.

The remedy follows from the stage objectives: since stage 1 optimises recall, it
uses gazetteer IDF alone, and only the feature representation uses the product.
Stage-1 recall rose from 81.8\% to 88.2\%.

This was invisible in aggregate recall, which moved by a few points and could
plausibly have been accepted. I found it by listing which tickers were being
missed and recognising the names.

\subsection{Defect 2: Compound-Name Fragmentation}
\label{sec:compound}

ASR renders CamelCase company names as separate words: CrowdStrike as ``Crowd
Strike'', SoFi as ``So Fi'', DraftKings as ``Draft Kings''. The fragments are
ordinary vocabulary, match nothing, and the company disappears from the pool.
Concatenating adjacent capitalised tokens and testing the result recovers these
cases, raising recall from 88.2\% to 88.9\%.

This is the same class of ASR damage as the phonetic corruption of ``Mitek''
into ``Midic'', differing only in that the error is a boundary error rather than
a substitution.

\subsection{Resulting Recall}

Adding dense retrieval over company-name embeddings and OCR symbols gives
\textbf{92.6\%} stage-1 recall, at a median of 93 and mean of 147 candidates per
segment. This figure is the hard ceiling on system performance and I report it
first in all subsequent analysis.

\section{Stage 2: Learned Ranking}
\label{sec:rank}

\subsection{Why Ranking, and Why Learned}

Since the correct symbol is in the pool 92.6\% of the time, segment attribution is
a ranking problem over a small set, not a generation problem over a large
vocabulary.

My initial scorer assigned weights by hand: 3.0 for symbol presence, 4.0 for
cue-phrase adjacency, and so on. It achieved recall@1 of 58.6\% against
recall@5 of 82.2\%. That gap is the diagnosis. The evidence was enough to put
the answer in the top five four times out of five, and the numbers I had chosen
were failing to put it first.

\begin{table}[htbp]
\caption{Hand-weighted scores for the Veritone segment. Correct answer is VERI.}
\label{tab:veri}
\centering
\begin{tabular}{rll}
\toprule
Score & Symbol & Company \\
\midrule
\textbf{12.30} & \textbf{VERI} & Veritone, Inc. \\
11.30 & V & Visa Inc. \\
4.80 & NATNF & Nation Gold Corp. \\
4.80 & LYV & Live Nation Entertainment \\
\bottomrule
\end{tabular}
\end{table}

Table~\ref{tab:veri} shows a representative near-failure. Visa is almost
selected, because the dictated spelling \texttt{V-E-R-I} contains the token
\texttt{V}, which is Visa's symbol, and the cue phrase ``ticker'' immediately
precedes it. Every rule that fires for VERI also fires for V. Telling them apart
needs the pieces of evidence weighed together, which fixed numbers chosen by
hand cannot do. The margin is 1.0 point.

With 687 labelled segments available, these weights are better estimated than
assumed.

\subsection{Feature Representation}

Each segment is treated as a query and its candidate pool as the ranked list,
giving 100{,}867 (segment, candidate) pairs over 37 features in six groups.

\emph{Symbol evidence}: presence, occurrence count, dictated form, cashtag,
cue-phrase adjacency, and normalised position of first mention. Position is
included because these segments announce the stock early; a mention in the final
fifth is usually a comparison rather than the subject.

\emph{Name evidence}: IDF-weighted mass of matched name tokens, exact token
count, coverage fraction, best fuzzy ratio, position of first mention.

\emph{Dense evidence}: cosine similarity and within-list rank between the
segment embedding and the company-name embedding, both from a MiniLM sentence
encoder~\cite{minilm}.

\emph{Provenance}: which stage-1 source contributed the candidate, and how many
did.

\emph{Structural alignment}: the segment's ordinal position within its video,
the candidate's first-mention position in the description and full transcript,
and the absolute difference between them. These videos are predominantly
countdown formats, so the $n$th segment typically concerns the $n$th enumerated
stock. No property of the segment's content encodes this.

\emph{Visual evidence}: Section~\ref{sec:ocr}.

\subsection{Model and Training}

I fit LambdaMART~\cite{lambdamart} using the LightGBM implementation.

A ranker is used rather than a classifier on purpose. A classifier asks whether
each candidate is correct on its own, which is the wrong question. The question
that matters is whether VERI is placed above V, and LambdaMART optimises that
ordering directly.

Training configuration: 900 boosting iterations, learning rate 0.03, 63 leaves,
minimum 15 samples per leaf, row and feature subsampling at 0.85, $L_2$
regularisation 1.0. Labels are binary, one relevant document per query. Each
fold is fitted three times under different random seeds and the predicted scores
averaged before selection.

I note one practical constraint discovered while testing a single-stage variant:
LightGBM rejects any query containing more than 10{,}000 documents. Ranking the
full registry per segment is therefore not expressible in this implementation
without truncation, independent of its cost.

\section{Visual Evidence}
\label{sec:ocr}

\subsection{Motivation}

The benchmark attributes the multimodal advantage on Task T specifically to
on-screen charts displaying ticker names. If that is the operative mechanism,
what is needed is reading text, not understanding images, and reading text is
far cheaper.

The hypothesis is directly testable: if symbols are written on screen, OCR of
sampled frames should recover them, and should do so disproportionately in
segments where the symbol is never spoken.

\subsection{Acquisition}

The dataset distributes source URLs rather than frames, so frames must be
fetched. I resolve a direct stream URL per video and extract six evenly spaced
JPEG frames per segment using HTTP range requests, so no video is downloaded in
full.

\subsection{Extraction}

Frames are processed with a GPU OCR model. I treat two patterns as strong
evidence because they are chart-header conventions rather than incidental text:
parenthesised symbols, \texttt{(VERI)}, and exchange-prefixed symbols,
\texttt{NASDAQ: VERI}. In practice the dominant source is a screen-recorded
financial quote page whose header reads, for example, \emph{``Veritone, Inc.
(VERI)''}.

Frames also contain substantial irrelevant text: bookshelves, wall decorations,
unrelated captions. This stage therefore makes no selection. It emits evidence
annotated with a confidence and a pattern type, which the ranker weighs against
everything else.

\begin{table}[htbp]
\caption{Visual evidence coverage}
\label{tab:ocr}
\centering
\begin{tabular}{lr}
\toprule
Quantity & Value \\
\midrule
Videos retrieved & 158 / 215 \\
Videos unavailable (deleted or private) & 57 (26.5\%) \\
Frames extracted & 2{,}628 \\
Segments with any frame & 438 / 687 (63.8\%) \\
\textbf{Correct symbol readable on screen} & \textbf{288 / 687 (41.9\%)} \\
\quad matching a chart-header pattern & 97 \\
\bottomrule
\end{tabular}
\end{table}

\subsection{Result}

The 41.9\% in Table~\ref{tab:ocr} is computed over all 687 segments, including
the 36.2\% for which no frame could be obtained; among segments with frames it
is 65.8\%.

For comparison the transcript contains the symbol in some literal form in
52.1\% of segments. The two
channels are thus comparably informative and, critically, informative in
\emph{different} segments: the screen continues to display \texttt{VERI} where
the presenter says only ``Veritone''. It is that difference, not the raw coverage, that makes the
video worth collecting.

\section{Experimental Protocol}

\subsection{Grouped Cross-Validation}

I use five-fold cross-validation with folds partitioned by \texttt{video\_id}
via \texttt{GroupKFold}. Five folds split 215 videos into test sets of roughly
43, large enough to hold a representative mix of video formats while still
leaving 80\% of the data for fitting. The grouping matters. Segments from a
common video share a candidate pool, a description, a title and a full
transcript; a
random partition would place near-duplicates on both sides and the resulting
estimate would be optimistic by an unknown margin. Under grouping no video
appears in both training and test. All reported figures are out-of-fold.

\subsection{Seed Averaging}

LightGBM subsamples rows and features stochastically, so a single seed yields a
non-deterministic estimate. On 687 queries I observe approximately one point of
F1 variation across seeds; one configuration scored 78.31 on a single seed and
78.02 averaged over three. Since the margins separating adjacent published
systems are of comparable size, a single-seed result cannot distinguish a
genuine improvement from sampling noise. All reported scores are averaged over
three seeds. Three is a pragmatic choice: it reduces the spread enough that the
remaining variation is well below the margins being claimed, at three times the
training cost rather than ten.

\section{Results}

\begin{table}[htbp]
\caption{Task T, segmented input. Published figures from~\cite{videoconviction},
Table 6.}
\label{tab:results}
\centering
\begin{tabular}{llr}
\toprule
System & Input & F1 \\
\midrule
Mistral-7B & text & 63.73 \\
DeepSeek-R1 & text & 71.18 \\
Gemini 1.5 Flash & text & 72.09 \\
Claude 3.5 Sonnet & text & 75.90 \\
Gemini 1.5 Pro & text & 76.56 \\
GPT-4o & text & 76.76 \\
DeepSeek-V3 & text & 77.56 \\
Llama-3.1-405B & text & 78.60 \\
Qwen2.5-72B & text & 79.20 \\
\midrule
\textbf{This work} & \textbf{text + frame OCR} & \textbf{80.06} \\
\midrule
Gemini 2.0 Flash & text + video & 82.24 \\
GPT-4o & text + video & 83.47 \\
Gemini 1.5 Pro & text + video & 86.01 \\
Gemini 1.5 Flash & text + video & 86.23 \\
\bottomrule
\end{tabular}
\end{table}

F1 80.06 exceeds all 16 text-only published results, including the strongest,
while using no large language model, no API and a single consumer GPU. The five
systems above it are all multimodal.

\subsection{Decomposition}

The score splits exactly into the two stages:
\begin{equation*}
  \mathrm{F1} \;=\; R_{1} \times A_{2} \;=\; 0.926 \times 0.865 \;=\; 0.801
\end{equation*}

where $R_{1}$ is stage-1 recall, the share of segments whose correct ticker
reached the candidate pool, and $A_{2}$ is ranker accuracy, the share of those
the ranker then placed first.

This holds exactly, not approximately. The system can only answer with a
candidate stage 1 produced, so a segment whose answer never entered the pool is
wrong no matter what the ranker does. Every segment receives exactly one answer
and has exactly one correct one, so precision, recall and F1 coincide, and the
fraction correct is the probability the answer reached the pool multiplied by
the probability the ranker then chose it.
Each term is independently improvable, and the factorisation says which to
target. This is the principal practical advantage of the decomposition over an
end-to-end method: an aggregate score does not distinguish failing to consider
the answer from failing to select it.

\subsection{Ablation}

\begin{table}[htbp]
\caption{Contribution of each component}
\label{tab:ablation}
\centering
\begin{tabular}{lrr}
\toprule
Configuration & Stage-1 recall & F1 \\
\midrule
Flat lexical, segment only & --- & 54.4 \\
Flat lexical + description & --- & 47.7 \\
Two-stage, hand-weighted scorer & 88.9\% & 56.6 \\
\quad + learned ranker & 88.6\% & \textbf{73.1} \\
\quad + dense similarity feature & 88.6\% & 75.1 \\
\quad + structural alignment & 88.6\% & 75.7 \\
\quad + dense retrieval in stage 1 & 91.6\% & 78.0 \\
\quad + frame OCR & \textbf{92.6\%} & \textbf{80.1} \\
\bottomrule
\end{tabular}
\end{table}

Table~\ref{tab:ablation} supports three observations. Adding the description
to a flat extractor is actively harmful, as established in
Section~\ref{sec:rationale}. The single largest gain,
16.5 points, comes from replacing hand-chosen weights with a learned ranker
rather than from any additional signal. Frame OCR contributes the final 2.0
points and is the step that crosses the strongest text-only baseline.

\subsection{Feature Importance}

The two dense features rank highest by split count, followed by candidate-pool
size and the stage-1 ordering. Lexical symbol features carry lower aggregate
importance but are not redundant: they are decisive on the 52.1\% of segments
where the symbol is present and silent elsewhere. The complementarity between
lexical and dense evidence reported in document retrieval~\cite{luan} therefore
reproduces here at the level of individual features.

\subsection{Selective Prediction}

A misidentified symbol is more damaging in any downstream analysis than a
declined one. The ranker's margin, the score gap between its first and second
candidate, provides an abstention signal requiring no additional model.

\begin{table}[htbp]
\caption{Selective prediction by ranking margin}
\label{tab:abstain}
\centering
\begin{tabular}{lrr}
\toprule
Operating point & Coverage & Precision \\
\midrule
Answer all & 100.0\% & 80.06 \\
Drop lowest 10\% by margin & 90.0\% & 86.25 \\
Drop lowest 20\% & 79.9\% & 90.71 \\
Drop lowest 30\% & 70.0\% & \textbf{93.76} \\
\bottomrule
\end{tabular}
\end{table}

Precision rises steadily across the operating points in
Table~\ref{tab:abstain}, which means the margin says something real about
whether the answer is right, not just about how many candidates there were.

\section{Limitations}

\textbf{Comparability.} The published systems are evaluated zero-shot; mine is
trained on VideoConviction under grouped cross-validation. The result is valid
as a benchmark figure on this dataset and is not a claim about transfer to a
different corpus. I state this prominently because Table~\ref{tab:results} would
otherwise be misread.

\textbf{Corpus attrition.} 26.5\% of source videos are no longer retrievable, so
36.2\% of segments receive no visual evidence. The benchmark authors conducted
their experiments when a larger fraction of the corpus was available, so their
multimodal figures reflect fuller visual coverage than is now obtainable.

The size of that effect can be estimated from the system's own predictions.
Partitioning the 687 segments by whether any frame was retrieved gives F1
$84.47$ on the 438 segments that still have video against $72.29$ on the 249
that do not, a gap of $12.19$ points. The 438-segment figure is within two
points of the strongest multimodal baseline, which had full video access. This
is not a clean experiment. Segments whose video survives carry transcripts
46\% longer on average, 2{,}342 characters against 1{,}605, which independently
favours the text features, so the split is an upper bound on what restored
coverage would recover rather than a prediction. A gap of this size is
nonetheless difficult to attribute to method alone, and it places most of the
distance between 80.06 and 86.23 on corpus attrition.

\textbf{Partial visual signal.} The multimodal baselines process video. This
system reads text from still frames and uses no motion, layout or chart
geometry.

\textbf{Single corpus.} As with any single-dataset result, the method cannot be
fully separated from the properties this corpus rewards.

\textbf{Label noise.} The answer key contains entries such as
\texttt{"NYSE: AMZN"} and \texttt{"ADR"}, and approximately 24\% of the correct answers
no longer trade.

\section{Conclusion}

Treating ticker extraction as a generation problem hides the fact that the
answer is often missing from the input, and missing more often for bigger
companies. Reformulating the task as retrieval followed by ranking makes the binding
constraint explicit and separately measurable, and reading symbols from video
frames with OCR captures a substantial part of the multimodal advantage at
negligible cost. The resulting system exceeds every text-only published baseline
without a large language model.

Two directions follow directly from the decomposition. Stage-1 recall is limited
chiefly by colloquial-versus-registered company names and by single-character
symbols, suggesting alias resolution. Ranker accuracy is limited by a bi-encoder
representation that scores segment and candidate independently, suggesting a
cross-encoder.

A third direction interests me more. Everything measured here concerns
\emph{accidental} obfuscation: ASR corruption and the ordinary pragmatics of not
naming a familiar company. A finfluencer subject to regulatory attention has a
direct incentive to avoid automated attribution, which is the structure of an
adversary rewriting content to evade a detector under a query budget. Whether
any of the signals used here survive that setting is an open question, and not
one the present evaluation addresses.

\section*{Code and Data}

The dataset is public. All preparation, analysis and evaluation code described
here, together with the hand-built symbol registry, is available from the
author. Every figure in this paper is produced by one of four scripts and
written to a results file, so each number can be traced to the run that
generated it.

\begin{thebibliography}{9}

\bibitem{videoconviction}
M. Galarnyk, V. Kejriwal, A. Shah, Y. Bhardwaj, N. W. Meyer, A. Krishnan, and
S. Chava, ``VideoConviction: A multimodal benchmark for human conviction and
stock market recommendations,'' in \emph{Proc. ACM SIGKDD Conf. Knowledge
Discovery and Data Mining (KDD)}, Toronto, ON, Canada, 2025.

\bibitem{lambdamart}
C. J. C. Burges, ``From RankNet to LambdaRank to LambdaMART: An overview,''
Microsoft Research, Tech. Rep. MSR-TR-2010-82, 2010.

\bibitem{luan}
Y. Luan, J. Eisenstein, K. Toutanova, and M. Collins, ``Sparse, dense, and
attentional representations for text retrieval,'' \emph{Trans. Assoc. Comput.
Linguistics}, vol. 9, pp. 329--345, 2021.

\bibitem{minilm}
N. Reimers and I. Gurevych, ``Sentence-BERT: Sentence embeddings using Siamese
BERT-networks,'' in \emph{Proc. EMNLP-IJCNLP}, Hong Kong, 2019, pp. 3982--3992.

\end{thebibliography}

\end{document}
